% frontmatter/lupa/mapa-del-curso.tex -- el año entero en una tabla.
\section*{Mapa del curso}

\begin{center}
\renewcommand{\arraystretch}{1.4}
\begin{tabularx}{\linewidth}{@{}c c c c X@{}}
\toprule
\textbf{Trimestre} & \textbf{Días} & \textbf{Época} & \textbf{Palabras} & \textbf{El gran caso} \\
\midrule
1 & 1--65   & otoño & 55--85 & Llega Hugo, nace el Club de la Lupa, y unas
notas misteriosas firmadas por \textit{Lector X} aparecen por casa.
¿Quién las escribe? \\
2 & 66--130 & invierno & 85--100 & En Nochebuena, Lucía encuentra el mapa
del tesoro que la abuela Rosa dibujó cuando tenía ocho años. \\
3 & 131--195 & primavera & 100--115 & Hace veinticinco años, una clase
enterró una cápsula del tiempo en el patio del colegio. Nadie recuerda
dónde. \\
4 & 196--260 & verano & 110--125 & En el pueblo, con el mapa de la
abuela: ¿dónde está el tesoro de Rosa? \\
\bottomrule
\end{tabularx}
\end{center}

\vspace{6mm}
Cada trimestre tiene 13 semanas de 5 días, también durante las
vacaciones. Lo que crece durante el año no es solo la longitud del texto
(la columna \textit{Palabras}): también lo largas que son las oraciones
y cuántas partes tienen, el tiempo en que se cuenta la historia --
presente en otoño, pasado desde el invierno -- y el tipo de texto:
notas, listas, cartas, recetas, instrucciones, entrevistas, textos
informativos, un diario, poemas y noticias.
\newpage
